\chapter{Introdução}
\label{cap:introducao}

A crescente complexidade da gestão urbana moderna exige que a Administração Pública adote infraestruturas tecnológicas capazes de processar, organizar e democratizar o acesso a grandes volumes de dados geoespaciais do município. Nesse cenário, os Sistemas de Informação Geográfica (SIG), do inglês \textit{Geographic Information System} (GIS), baseados na \textit{web} (WebGIS), consolidaram-se como ferramentas essenciais para o planejamento urbano e a transparência governamental. O uso dessas ferramentas tem se expandido no setor público \cite{Silva2013} como suporte fundamental para a tomada de decisões \cite{Alcantara2015}. Acompanhando essa expansão, a arquitetura desses sistemas tem evoluido rapidamente, impulsionando pesquisas focadas em escalabilidade, interoperabilidade e alto desempenho \citeonline{VinuezaMartinez2024}.

Apesar da importância comprovada dessas ferramentas, o desenvolvimento de aplicações cartográficas para a \textit{web} apresenta desafios próprios da engenharia de \textit{software}. A publicação de mapas topográficos e dados vetoriais exige métodos eficientes de renderização e transferência de rede para evitar a degradação da experiência do usuário \cite{BastosCamboim2022}. Quando o volume de camadas (\textit{layers}) e o tráfego de usuários aumentam, aplicações estruturadas sob modelos arquiteturais monolíticos tradicionais, ou dependentes de \textit{frameworks} legados, tendem a apresentar lentidão, interfaces pouco responsivas para dispositivos móveis e gargalos significativos de desempenho e manutenção \cite{Marotta2020}.

É neste contexto de inovação e modernização tecnológica que se consolida o projeto "Sobral em Mapas", uma iniciativa estratégica e uma ferramenta imprescindível para a Prefeitura Municipal de Sobral, que busca integrar e publicar os dados geoespaciais do município utilizando a biblioteca OpenLayers. Durante o desenvolvimento do projeto, identificou-se que, para assegurar a longevidade, a segurança e a escalabilidade da plataforma em um ambiente de produção, seria necessário adotar um ecossistema tecnológico robusto, superando as limitações de servidores locais e pilhas de desenvolvimento convencionais.

O presente trabalho aborda essa evolução por meio da transição para uma arquitetura de \textit{software} estritamente desacoplada. Ao separar a interface do usuário (\textit{front-end}), desenvolvida como uma \textit{Single Page Application} (SPA) reativa utilizando a biblioteca React, da lógica de controle e acesso a dados (\textit{back-end}), orquestrada por uma API \textit{RESTful} em PHP Laravel, assim o sistema mitiga gargalos de recarregamento e ganha flexibilidade. Além disso, para solucionar as complexidades históricas de implantação (\textit{deploy}) e garantir a reprodutibilidade do ambiente em qualquer infraestrutura, a aplicação foi integralmente conteinerizada com a tecnologia Docker. Em paralelo, o servidor GeoServer foi adotado como o motor robusto para a distribuição inteligente de dados abertos padronizados.

\section{Justificativa}
A justificativa para este trabalho fundamenta-se na urgência de modernizar a gestão de dados geográficos do município de Sobral, superando as limitações de sistemas monolíticos legados por meio da consolidação do projeto "Sobral em Mapas". A transição para uma arquitetura estritamente desacoplada e conteinerizada resolve gargalos críticos de desempenho e implantação (\textit{deploy}). Integrada a essa base tecnológica, a adoção da biblioteca \textit{open-source} OpenLayers fornece um motor de renderização flexível e compatível com padrões interoperáveis (como WMS e WFS), viabilizando a exploração ágil do espaço urbano para embasar a formulação de políticas públicas e a tomada de decisão governamental.

Para além dos avanços computacionais, o projeto cumpre um papel estratégico ao democratizar o acesso ao conhecimento cartográfico local, disponibilizando informações espaciais de forma ampla à sociedade \cite{Marotta2020} e consolidando o WebGIS como tecnologia ascendente na representação territorial \cite{Alcantara2015}. Simultaneamente, a construção deste ecossistema atua como um laboratório de validação científica que favorece o aprimoramento em geotecnologias e programação \textit{web}, firmando-se como um recurso essencial não apenas para o planejamento urbano de Sobral, mas também para a formação acadêmica qualificada.

\subsection{Objetivo Geral}
O objetivo geral deste trabalho é projetar, desenvolver e implementar uma arquitetura de \textit{software} moderna, estritamente desacoplada e conteinerizada para a plataforma \textit{web} do projeto "Sobral em Mapas", buscando construir um ecossistema tecnológico escalável e de alto desempenho que otimize a gestão, a integração, a distribuição e a visualização interativa de dados geoespaciais abertos, fornecendo assim uma ferramenta resiliente e de fácil manutenção para o apoio à gestão pública do município de Sobral.

\subsection{Objetivos Específicos}
Para que o objetivo geral seja alcançado com êxito, a metodologia deste trabalho foi estruturada a partir dos seguintes objetivos específicos:
\begin{itemize}
    \item Realizar o levantamento de requisitos e a modelagem estrutural do sistema, definindo as entidades de banco de dados e os relacionamentos necessários para o armazenamento de usuários, categorias e metadados geográficos;
    \item Projetar e desenvolver o \textit{back-end} da aplicação sob o formato de uma API \textit{RESTful} utilizando o \textit{framework} PHP Laravel, implementando rotas de comunicação, autenticação segura e orquestrando toda a lógica de negócios da plataforma;
    \item Construir a interface com o usuário (\textit{front-end}) adotando o paradigma de \textit{Single Page Application} (SPA) por meio da biblioteca React, com foco na componentização do código e na entrega de uma experiência de usuário (UX) fluida e altamente responsiva em múltiplos dispositivos;
    \item Integrar a infraestrutura de dados ao servidor de mapas de código aberto GeoServer para a publicação e administração inteligente das bases espaciais, utilizando a biblioteca OpenLayers no cliente para o consumo, manipulação interativa e renderização em tempo real das camadas cartográficas;
    \item Padronizar os ambientes de desenvolvimento e produção por meio da conteinerização com a tecnologia Docker, estabelecendo contêineres independentes para os serviços da aplicação, o que garante um fluxo de implantação (\textit{deploy}) isolado, reprodutível e altamente escalável.
\end{itemize}